\documentclass[11pt,a4paper]{article}

\usepackage{fontspec}
\usepackage[main=russian, english]{babel}
\usepackage{amsmath, amsfonts, amssymb}
\usepackage{graphicx}
\usepackage{booktabs}
\usepackage{cite}
\usepackage{geometry}
\usepackage{xcolor}
\usepackage{url}
\usepackage{hyperref}

% Настройка шрифтов для XeLaTeX/LuaLaTeX (по умолчанию используются системные)
\setmainfont{Times New Roman}
\setromanfont{Times New Roman}
\setsansfont{Arial}
\setmonofont{Courier New}

\geometry{margin=2.5cm}

\title{\textbf{Адаптивное извлечение секвенограмм: слепая деконволюция спектрограмм с помощью физически обоснованной шаблонной ConvNMF}}
\author{Ваше Имя\thanks{Автор для корреспонденции: email@example.com} \\ 
        \small \textit{Кафедра вычислительной обработки сигналов}}
\date{\today}

\begin{document}

\maketitle

\begin{abstract}
В сложных акустических условиях многолучевое распространение и реверберация значительно искажают спектрограммы сигналов, маскируя временную динамику звуковых эмиссий. Мы представляем новый вычислительный подход для извлечения \textit{секвенограмм} — высокоразрешенных, разреженных векторов временной активации — из реверберирующих спектрограмм с использованием словаря заданных шаблонов. Наш метод использует физически обоснованный подход сверточной неотрицательной матричной факторизации (ConvNMF). Вместо использования стандартных 2D-ядер, мы моделируем экологическое эхо с помощью частотно-зависимого причинного 1D-фильтра. Это естественным образом учитывает физическое явление частотно-зависимого поглощения звука и предотвращает структурные искажения эталонных шаблонов в процессе оптимизации. Валидация на синтетических данных демонстрирует высокую надежность алгоритма при слепой оценке профиля эха и локализации шаблонов. Предложенный метод представляет собой вычислительно эффективный и адаптивный инструмент для поиска паттернов в задачах крупномасштабного акустического мониторинга.
\end{abstract}

\section{Введение}
Точная локализация и классификация сигнальных паттернов на частотно-временных представлениях являются фундаментальными задачами в таких областях, как биоакустика, промышленный мониторинг и гидроакустика. В естественных средах (пещеры, леса или замкнутые городские пространства) за первичными сигналами неизбежно следует серия отражений. Эти отражения формируют «хвост эха», который размазывает акустическую энергию по временной оси, эффективно размывая спектрограмму.

Традиционные методы обнаружения сигналов, такие как взаимная корреляция или стандартные сверточные нейронные сети (CNN) для обработки изображений, крайне чувствительны к таким искажениям. Когда эхо одного сигнала накладывается на излучение последующего, стандартные алгоритмы часто демонстрируют высокий уровень ложных срабатываний и потерю временной точности.

Для решения этой проблемы мы вводим понятие \textbf{секвенограммы}: 1D разреженного вектора активаций, который представляет точное время излучения конкретных шаблонов, очищенное от артефактов окружающей среды. Наш подход использует метод физически обоснованной слепой деконволюции для извлечения этих секвенограмм с одновременным обучением адаптивного профиля эха среды записи.

\section{Материалы и методы}

\subsection{Физическая и генеративная модель}
Пусть $S \in \mathbb{R}_{\ge 0}^{F \times T}$ — наблюдаемая амплитудная или энергетическая спектрограмма, где $F$ — количество частотных бинов, а $T$ — количество временных кадров. Мы предполагаем, что генеративный процесс сигнала состоит из двух этапов: излучения и реверберации.

Сначала генерируется «чистая» (безэховая) спектрограмма $V$. Задан словарь из $K$ известных шаблонов $W_k$, каждый из которых соответствует определенному классу сигналов. Чистая спектрограмма формируется путем свертки этих шаблонов с их соответствующими неотрицательными векторами активации (секвенограммами) $H_k \in \mathbb{R}_{\ge 0}^{T}$:
\begin{equation}
V(f, t) = \sum_{k=1}^{K} (H_k *_{t} W_k(f, \cdot))(t)
\end{equation}
где $*_{t}$ обозначает свертку по временной оси.

Затем чистая спектрограмма $V$ подвергается реверберации среды. Вместо использования стандартной 2D-свертки, которая может ошибочно переносить энергию между частотными полосами, мы применяем частотно-зависимую 1D-матрицу эха $E \in \mathbb{R}_{\ge 0}^{F \times L_e}$. Каждая строка $E$ действует как независимая причинная импульсная характеристика для конкретного частотного канала:
\begin{equation}
\hat{S}(f, t) = (V(f, \cdot) *_{t} E(f, \cdot))(t)
\end{equation}
Такая архитектура естественно отражает физику распространения звука: высокочастотные звуковые волны затухают в воздухе быстрее, чем низкочастотные.

\subsection{Физически обоснованная калибровочная нормировка функционалов ошибки}
Для извлечения идеальных секвенограмм мы должны наложить структурные ограничения на векторы активации $f(t) \in [0,1]$ (где $f(t)$ — произвольный канал матрицы $H$). В частности, мы требуем, чтобы $f(t)$ состоял из изолированных бинаризованных пиков (событий дираковского типа) с ожидаемой физической плотностью $\rho_0$ (пиков в секунду).

В стандартном многозадачном обучении (Multi-Task Learning) структурные штрафы объединяются с помощью линейной комбинации $\mathcal{L}_{total} = \sum \alpha_i \mathcal{L}_i$, где $\alpha_i$ — настраиваемые вручную гиперпараметры. Линейный подход содержит фундаментальный геометрический недостаток: градиенты действуют как плоская поверхность, непрерывно применяя возвращающую силу даже тогда, когда конкретный компонент ошибки достигает нуля, что может выталкивать сеть в отрицательные, нестабильные режимы.

Вместо этого мы предлагаем агрегировать компоненты ошибки, используя \textbf{Евклидову метрику} $\mathcal{L}_{total} = \sqrt{\sum \mathcal{L}_i^2}$. Евклидова геометрия создает коническую потенциальную яму в нулевой конфигурации. По мере приближения компонента $\mathcal{L}_i$ к нулю, вклад его частной производной естественно исчезает, динамически переключая фокус оптимизации на оставшиеся невыполненные ограничения.

Однако евклидов подход справедлив только в том случае, если все функциональные компоненты $\mathcal{L}_i[f]$ находятся в изотропном пространстве. Мы достигаем этого путем формулирования системы \textit{Калибровочной нормировки (Gauge Calibration)}. Любой допустимый функционал должен удовлетворять трем строгим физическим условиям:
\begin{enumerate}
    \item \textbf{Нулевое основное состояние:} $\mathcal{L}[f^*] = 0$ для ожидаемой идеальной конфигурации $f^*$.
    \item \textbf{Единичная вариационная возвращающая сила:} Средняя величина вариационной производной в ожидаемой конфигурации должна быть равна единице: $\langle | \frac{\delta \mathcal{L}}{\delta f(t)} | \rangle_{f^*} = 1$.
    \item \textbf{Инвариантность к дискретизации:} Функционал должен сходиться к стабильному непрерывному интегралу при увеличении частоты дискретизации ($\Delta t \to 0$).
\end{enumerate}

Основываясь на этих принципах, мы выводим два основных штрафных функционала для секвенограмм. Пусть полное физическое время равно $\tau_{max}$, а ожидаемая длительность события — $\tau_e$ (как правило, предел одного дискретного отсчета, $\tau_e = \Delta t$).

\subsubsection*{1. Функционал резкости (Бинаризация)}
Мы штрафуем отклонения от значений $0$ и $1$. Непрерывный функционал определяется как:
\begin{equation}
\mathcal{L}_{sharp}[f] = \int_{0}^{\tau_{max}} f(t)\big(1 - f(t)\big) dt
\end{equation}
\textit{Доказательство выполнения условий:}
1) Для $f(t) \in \{0, 1\}$ подынтегральное выражение строго равно нулю (выполняется условие основного состояния). 
2) Вариационная производная равна $\frac{\delta \mathcal{L}_{sharp}}{\delta f(t)} = 1 - 2f(t)$. Поскольку $f(t) \in \{0,1\}$, норма модуля $|1-2f(t)| = 1$ повсюду. Таким образом, ее среднее значение равно $1$.
3) Дискретная аппроксимация $\mathcal{L}_{sharp} \approx \Delta t \sum_i f_i(1-f_i)$ явно включает $\Delta t$, обеспечивая инвариантность при изменении частоты дискретизации.

\subsubsection*{2. Функционал изоляции (Оконный штраф)}
Для предотвращения размытия пиков и чрезмерно частых срабатываний мы применяем окно локального отталкивания. Пусть $w(t)$ — индикаторная функция, равная $1$ для $t \in [-\tau_w, \tau_w]$ и $0$ в остальных местах, с выколотым центром $w(0)=0$. Откалиброванный функционал имеет вид:
\begin{equation}
\mathcal{L}_{iso}[f] = \lambda \int_{0}^{\tau_{max}} f(t) \big(w * f\big)(t) dt
\end{equation}
\textit{Доказательство выполнения условий:}
1) Если пики разделены расстоянием, большим $\tau_w$, свертка равна $0$ в позиции каждого пика. Следовательно, $\mathcal{L}_{iso}[f^*] = 0$.
2) Вариационная производная равна $\frac{\delta \mathcal{L}_{iso}}{\delta f(t)} = 2\lambda (w * f)(t)$. В ожидаемой конфигурации, содержащей $N = \rho_0 \tau_{max}$ пиков шириной $\tau_e$, свертка равна $\tau_e$ внутри окна каждого пика. Ненулевая область занимает время $N \times 2\tau_w$. Средняя вариационная норма на интервале $\tau_{max}$ составляет:
\begin{equation}
\langle | \frac{\delta \mathcal{L}_{iso}}{\delta f} | \rangle = \frac{1}{\tau_{max}} \left( N \cdot 2\tau_w \right) \cdot (2\lambda \tau_e) = 4 \lambda \rho_0 \tau_w \tau_e
\end{equation}
Для обеспечения единичной вариационной возвращающей силы (Условие 2) мы однозначно определяем масштабную константу $\lambda$:
\begin{equation}
\lambda = \frac{1}{4 \rho_0 \tau_w \tau_e}
\end{equation}
3) В дискретном времени, если ширина окна составляет $W = \tau_w / \Delta t$ бинов, а ширина пика равна $1$ бину ($\tau_e = \Delta t$), дискретный откалиброванный функционал принимает вид:
\begin{equation}
\mathcal{L}_{iso} = \frac{\Delta t}{4 \rho_0 \tau_w} \sum_{i} f_i (w * f)_i
\end{equation}
Данная формулировка сохраняет инвариантность к дискретизации, автоматически адаптируясь к физическим параметрам среды.

\subsubsection*{3. Единая изотропная целевая функция}
Поскольку оба функционала откалиброваны так, чтобы проецировать единичную градиентную силу вблизи решения, они образуют идеально изотропное пространство ошибок. Мы объединяем их с использованием Евклидовой нормы:
\begin{equation}
\mathcal{L}_{reg} = \sqrt{\mathcal{L}_{sharp}^2 + \mathcal{L}_{iso}^2 + \epsilon}
\end{equation}
где $\epsilon$ — малая константа, обеспечивающая численную стабильность в нуле. Данная регуляризация добавляется к основной функции потерь реконструкции спектрограммы (KL-дивергенция).

\section{Результаты}

\subsection{Синтетическая валидация (Проверка концепции)}
Для проверки способности алгоритма разделять перекрывающиеся сигналы и адаптивное эхо был создан синтетический набор данных с известной истинной разметкой. Набор данных включал два шаблона, занимающих разные частотные полосы (низкую и высокую). Были сгенерированы разреженные пики активации для обоих шаблонов, после чего применялось частотно-зависимое экспоненциальное затухание в качестве эха. Наконец, для имитации реальных условий записи был добавлен гауссов шум.

Модель слепой деконволюции успешно реконструировала разреженные секвенограммы с точной временной привязкой, различая перекрывающиеся вызовы. Более того, изученная матрица $E$ точно соответствовала истинным скоростям затухания, автоматически выявив, что высокие частоты затухают быстрее низких.

% [TODO: Вставить Рисунок 1: Синтетическая спектрограмма, извлеченные секвенограммы и профили эха]

\subsection{Применение к реальным акустическим данным}
% [TODO: Описать результаты, полученные на спектрограмме эхолокации летучих мышей размером 600,000x300. Упомянуть, как выбирались шаблоны и как извлеченные секвенограммы улучшили точность временной привязки сигналов по сравнению с исходными данными.]

\section{Обсуждение}
Главным преимуществом предложенной частотно-зависимой 1D-свертки перед стандартной 2D-деконволюцией является ее физическая согласованность. Запрещая перенос энергии между частотами внутри модели эха, мы заставляем алгоритм объяснять временное размытие строго через реверберацию. Это предотвращает искусственное искажение формы эталонных шаблонов моделью ради минимизации функции потерь.

Кроме того, мы представили систему Калибровочной нормировки (Gauge Calibration) для структурных функций потерь. Математически обеспечив \textit{единичную вариационную возвращающую силу} для всех штрафных членов, мы отобразили процесс оптимизации в изотропное пространство. Это позволило нам агрегировать структурные ограничения с использованием евклидовой геометрии ($\sqrt{\sum \mathcal{L}_i^2}$) вместо традиционных линейных комбинаций. Мы полагаем, что если исследователи примут эти три физических правила (Нулевое основное состояние, Единичная вариационная сила и Инвариантность к дискретизации) для регуляризации нейронных сетей, сообщество машинного обучения сможет полностью избавиться от необходимости эвристического подбора гиперпараметров при балансировке многозадачных целевых функций.

\section{Заключение}
Платформа адаптивного извлечения секвенограмм предлагает надежный, физически обоснованный подход к поиску паттернов в условиях сильной реверберации. Разделяя внутренний шаблон сигнала и внешнее эхо среды, метод обеспечивает высокоточные временные активации. Использование FFT-сверток и теоретически откалиброванных, масштабно-инвариантных функционалов потерь делает этот метод вычислительно эффективным и свободным от гиперпараметров инструментом для крупномасштабного анализа непрерывных акустических записей.

\section*{Доступность данных и кода}
Исходный код на Python (с использованием PyTorch) и скрипты для генерации синтетических данных доступны в нашем открытом репозитории: \url{https://github.com/yourusername/AdaptiveSequenograms}.

\bibliographystyle{plain}
\bibliography{references}

\end{document}